% Options for packages loaded elsewhere
\PassOptionsToPackage{unicode}{hyperref}
\PassOptionsToPackage{hyphens}{url}
\PassOptionsToPackage{dvipsnames,svgnames,x11names}{xcolor}
\documentclass[
  11pt,
  a4paper,
]{article}
\usepackage{xcolor}
\usepackage[margin=2.6cm]{geometry}
\usepackage{amsmath,amssymb}
\setcounter{secnumdepth}{5}
\usepackage{iftex}
\ifPDFTeX
  \usepackage[T1]{fontenc}
  \usepackage[utf8]{inputenc}
  \usepackage{textcomp} % provide euro and other symbols
\else % if luatex or xetex
  \usepackage{unicode-math} % this also loads fontspec
  \defaultfontfeatures{Scale=MatchLowercase}
  \defaultfontfeatures[\rmfamily]{Ligatures=TeX,Scale=1}
\fi
\usepackage{lmodern}
\ifPDFTeX\else
  % xetex/luatex font selection
  \setmainfont[]{Libertinus Serif}
  \setmonofont[Scale=0.85]{Latin Modern Mono}
  \setmathfont[]{Latin Modern Math}
\fi
% Use upquote if available, for straight quotes in verbatim environments
\IfFileExists{upquote.sty}{\usepackage{upquote}}{}
\IfFileExists{microtype.sty}{% use microtype if available
  \usepackage[]{microtype}
  \UseMicrotypeSet[protrusion]{basicmath} % disable protrusion for tt fonts
}{}
\makeatletter
\@ifundefined{KOMAClassName}{% if non-KOMA class
  \IfFileExists{parskip.sty}{%
    \usepackage{parskip}
  }{% else
    \setlength{\parindent}{0pt}
    \setlength{\parskip}{6pt plus 2pt minus 1pt}}
}{% if KOMA class
  \KOMAoptions{parskip=half}}
\makeatother
\usepackage{longtable,booktabs,array}
\newcounter{none} % for unnumbered tables
\usepackage{calc} % for calculating minipage widths
% Correct order of tables after \paragraph or \subparagraph
\usepackage{etoolbox}
\makeatletter
\patchcmd\longtable{\par}{\if@noskipsec\mbox{}\fi\par}{}{}
\makeatother
% Allow footnotes in longtable head/foot
\IfFileExists{footnotehyper.sty}{\usepackage{footnotehyper}}{\usepackage{footnote}}
\makesavenoteenv{longtable}
\setlength{\emergencystretch}{3em} % prevent overfull lines
\providecommand{\tightlist}{%
  \setlength{\itemsep}{0pt}\setlength{\parskip}{0pt}}

\providecommand{\E}{\mathbb{E}}
\providecommand{\N}{\mathbb{N}}
\providecommand{\Z}{\mathbb{Z}}
\providecommand{\R}{\mathbb{R}}
\providecommand{\eps}{\varepsilon}
\providecommand{\ceil}[1]{\left\lceil #1\right\rceil}
\providecommand{\floor}[1]{\left\lfloor #1\right\rfloor}
\AtBeginDocument{\let\setminus\smallsetminus}
\usepackage[most]{tcolorbox}
\newtcolorbox{warning}{colback=red!5,colframe=red!65!black,boxrule=0.8pt,arc=2pt,left=6pt,right=6pt,top=5pt,bottom=5pt}
\usepackage{etoolbox}
\AtBeginEnvironment{longtable}{\small}
\setlength{\emergencystretch}{3em}
\usepackage{fancyhdr}
\pagestyle{fancy}
\fancyhf{}
\fancyhead[L]{\small\itshape Candidate result 2008.03587\_\_00 --- GPT-6 Astra, awaiting review}
\fancyhead[R]{\small\thepage}
\renewcommand{\headrulewidth}{0.2pt}
\usepackage{bookmark}
\IfFileExists{xurl.sty}{\usepackage{xurl}}{} % add URL line breaks if available
\urlstyle{same}
\hypersetup{
  pdftitle={A candidate counterexample to Question 4.1 in A note on deterministic zombies},
  pdfauthor={github.com/graph-theory-AI --- machine-generated candidate writeup},
  colorlinks=true,
  linkcolor={blue!50!black},
  filecolor={Maroon},
  citecolor={Blue},
  urlcolor={blue!50!black},
  pdfcreator={LaTeX via pandoc}}

\title{A candidate counterexample to Question 4.1 in A note on
deterministic zombies}
\author{\href{https://github.com/graph-theory-AI}{github.com/graph-theory-AI}
--- machine-generated candidate writeup}
\date{17 September 2026}

\begin{document}
\maketitle
\begin{abstract}
A 59-vertex cactus requires at least three zombies, but adding one leaf
produces a graph requiring exactly two. This candidate disproof was
generated by GPT-6 Astra in a single model pass for the Graph Theory AI
project. It was selected from the campaign because the model marked it
\texttt{would\_publish:\ true}; it has not been checked by the
adversarial referee, a human mathematician, or a novelty search.
\end{abstract}

\section{The problem and the claimed
result}\label{the-problem-and-the-claimed-result}

This document concerns
\href{https://graph-theory-ai.github.io/graph-conjectures/arxiv/2008.03587__00/}{Question
4.1} from \emph{A note on deterministic zombies}.

\subsection{Problem statement}\label{problem-statement}

Question. For every graph \(G\), and for a graph \(G'\) obtained from
\(G\) by successively adding vertices of degree 1, does it always hold
that \(z(G') = z(G)\)?

\subsection{Catalog context}\label{catalog-context}

This question asks whether attaching pendant vertices (leaves) to a
graph leaves its zombie number unchanged. It can be interpreted as
asking whether zombies can benefit from waiting at a designated position
before activating at the start of the game.

\subsection{Claimed result}\label{claimed-result}

\textbf{Model verdict: disproved} (confidence: high).

A 59-vertex cactus requires at least three zombies, but adding one leaf
produces a graph requiring exactly two.

\emph{The model's own caveat:} Uses the standard compulsory-move rules,
with zombies moving before the survivor; no computational or
literature-status claim is made.

\section{Model writeup}\label{model-writeup}

The text below is the model output verbatim, apart from moving its
machine-readable verdict block into the summary above and shifting
Markdown heading levels for this document. Mechanical TeX defects and
equation tags were normalized where needed for compilation.

\subsection{A counterexample}\label{a-counterexample}

Zombies choose their initial positions first, followed by the survivor.
On each turn, every zombie must move along an edge decreasing its
distance to the survivor; the survivor then moves to a neighbor or
stays. Capture occurs upon coincidence.

The counterexample is a cactus consisting of three cycles attached to a
path. The proof below is entirely analytic.

\subsubsection{1. Construction}\label{construction}

Start with the path \[
v_0v_1\cdots v_{10}.
\] Attach, using otherwise disjoint vertices,

\begin{itemize}
\tightlist
\item
  a cycle \(A\) of length \(25\) at \(v_0\);
\item
  a cycle \(B\) of length \(17\) at \(v_2\);
\item
  a cycle \(C\) of length \(9\) at \(v_4\).
\end{itemize}

There are no other edges. Call the resulting graph \(G\). Thus \[
|V(G)|=11+24+16+8=59.
\]

Let \(H\) be obtained by adding a single new vertex \(v_{11}\), adjacent
only to \(v_{10}\). I will prove \[
z(G)\ge 3
\qquad\text{and}\qquad
z(H)=2.
\]

Here and below, each cycle includes its attachment vertex.

\begin{center}\rule{0.5\linewidth}{0.5pt}\end{center}

\subsubsection{2. A certificate for perpetual rotation around an odd
cycle}\label{a-certificate-for-perpetual-rotation-around-an-odd-cycle}

Let \(Q\) contain a cycle of length \[
n=2h+1,\qquad h\ge2,
\] attached to the rest of \(Q\) only at its vertex \(0\). Choose an
orientation and label its vertices by \(\mathbb Z_n\) in that
orientation.

For a zombie's initial position \(x\), define its \textbf{phase} \[
p(x)=
\begin{cases}
x, & x\text{ lies on the cycle, using its cyclic coordinate};\\[2mm]
-d_Q(x,0)\pmod n, & x\text{ lies outside the cycle}.
\end{cases}
\]

\paragraph{Rotation lemma}\label{rotation-lemma}

If there is a coordinate \(q\) such that \[
q-p(x_i)\pmod n\in\{2,3,\ldots,h\}
\qquad\text{(1)}
\] for every zombie \(i\), then the survivor can evade forever by
starting at \(q\) and moving one step in the chosen orientation every
turn.

In particular, two zombies can be evaded this way whenever \[
\|p(x_1)-p(x_2)\|_n\le h-2,
\qquad\text{(2)}
\] where \(\|a\|_n\) denotes cyclic distance to \(0\) modulo \(n\).

\paragraph{Proof}\label{proof}

First consider a zombie initially on the cycle. Condition (1) puts it a
forward distance between \(2\) and \(h\) behind the survivor. Its unique
shortest route is in the survivor's direction. After the zombie and
survivor each move, this forward distance is unchanged.

Now consider a zombie outside the cycle, initially at distance \(D\)
from the attachment vertex. While the survivor stays on the cycle, every
shortest route from that zombie to the survivor passes through \(0\).
Consequently, the zombie decreases its distance to \(0\) on every move
and reaches \(0\) on its \(D\)-th move.

Immediately before that landing move, the survivor is at \[
q+D-1 \pmod n.
\] Writing \(e=q+D\pmod n\), condition (1) gives \(2\le e\le h\). Thus
the survivor is not at \(0\) when the zombie lands. After the survivor's
response, the forward gap is \(e\), again between \(2\) and \(h\), and
remains constant thereafter.

For the final assertion, two residues can satisfy (1) simultaneously
exactly when their difference has a representative in \[
[-(h-2),h-2],
\] because this is the set of differences of two integers in
\(\{2,\ldots,h\}\). \(\square\)

For reference, the sufficient bounds in (2) are \[
\begin{array}{c|c|c}
\text{Cycle length}&h&h-2\\ \hline
25&12&10\\
17&8&6\\
9&4&2
\end{array}
\]

\begin{center}\rule{0.5\linewidth}{0.5pt}\end{center}

\subsubsection{\texorpdfstring{3. Two zombies do not suffice on
\(G\)}{3. Two zombies do not suffice on G}}\label{two-zombies-do-not-suffice-on-g}

Consider any two initial zombie positions \(x,y\).

Every vertex outside \(A\) has distance at most \(10\) from \(v_0\).
Indeed:

\begin{itemize}
\tightlist
\item
  the path ends at \(v_{10}\);
\item
  a vertex of \(B\) is at distance at most \(2+8=10\) from \(v_0\);
\item
  a vertex of \(C\) is at distance at most \(4+4=8\) from \(v_0\).
\end{itemize}

We distinguish all three possibilities for membership in \(A\).

\paragraph{\texorpdfstring{Case 1: Both zombies start outside
\(A\)}{Case 1: Both zombies start outside A}}\label{case-1-both-zombies-start-outside-a}

Put \[
D_x=d_G(x,v_0),\qquad D_y=d_G(y,v_0).
\] Both numbers lie in \(\{1,\ldots,10\}\), so \[
|D_x-D_y|\le9.
\]

Their phases on \(A\) are \(-D_x,-D_y\) modulo \(25\). Their cyclic
separation is therefore at most \(9\), which is below the rotation
lemma's bound \(10\). The survivor can rotate around \(A\) forever.

\paragraph{\texorpdfstring{Case 2: Both zombies start in
\(A\)}{Case 2: Both zombies start in A}}\label{case-2-both-zombies-start-in-a}

Let \[
a=d_G(x,v_0),\qquad b=d_G(y,v_0),\qquad
\delta=|a-b|.
\] Since \(A\) has length \(25\), \[
0\le\delta\le12.
\]

The zombies' distances to \(v_2\) are \(a+2,b+2\), and their distances
to \(v_4\) are \(a+4,b+4\). Thus their phase separation is represented
by \(\delta\) on both \(B\) and \(C\).

If \[
\delta\in\{0,1,\ldots,6,11,12\},
\] then \[
\|\delta\|_{17}\le6,
\] so the survivor can rotate around \(B\).

Otherwise \(\delta\in\{7,8,9,10\}\), and then \[
\|\delta\|_9\le2.
\] The survivor can instead rotate around \(C\).

\paragraph{\texorpdfstring{Case 3: Exactly one zombie starts in
\(A\)}{Case 3: Exactly one zombie starts in A}}\label{case-3-exactly-one-zombie-starts-in-a}

Suppose \(x\in A\) and \(y\notin A\). Put \[
t=d_G(x,v_0)\in\{0,\ldots,12\},
\qquad
D=d_G(y,v_0)\in\{1,\ldots,10\}.
\]

Orient \(A\) so that \(x\) has coordinate \(-t\pmod{25}\). The two
phases on \(A\) are then \(-t,-D\), whose difference is represented by
\[
D-t\in[-11,10].
\]

Except when \[
(D,t)=(1,12),
\] we have \(|D-t|\le10\), and the rotation lemma gives perpetual
evasion on \(A\).

In the exceptional case, \(y=v_1\), since this is the only vertex
outside \(A\) at distance \(1\) from \(v_0\). The distances of the
zombies to the attachment vertex \(v_2\) of \(B\) are \[
d_G(x,v_2)=12+2=14,
\qquad
d_G(y,v_2)=1.
\] Their phase difference on \(B\) is \(13\) modulo \(17\), with cyclic
distance \[
\|13\|_{17}=4\le6.
\] The survivor can rotate around \(B\).

These cases cover every pair of initial positions. Therefore \[
\boxed{z(G)\ge3.}
\]

\begin{center}\rule{0.5\linewidth}{0.5pt}\end{center}

\subsubsection{4. A delayed pair that captures on a pendant
cycle}\label{a-delayed-pair-that-captures-on-a-pendant-cycle}

The following observation supplies the winning strategy on \(H\).

\paragraph{Capture lemma}\label{capture-lemma}

Suppose a cycle of length \(2d+3\), with \(d\ge1\), is attached to the
rest of a graph only at its root \(r\). Suppose, immediately before a
zombie turn,

\begin{itemize}
\tightlist
\item
  one zombie is at \(r\);
\item
  another zombie is outside the cycle at distance \(d\) from \(r\);
\item
  the survivor is on the cycle.
\end{itemize}

Then these two zombies can force capture.

\paragraph{Proof}\label{proof-1}

Label the root \(0\). Choose the orientation in which the survivor's
initial coordinate \(p\) satisfies \[
1\le p\le d+1.
\] If \(p=1\), the root zombie captures immediately. Assume otherwise.

The root zombie follows the survivor in this orientation. Consider a
play avoiding capture, including capture on the following zombie turn.
After \(t\) complete rounds, write the first zombie's coordinate as
\(t\) and the survivor's unwrapped coordinate as \(u_t\).

The forward gap \[
g_t=u_t-t
\] cannot increase: on a survivor move it changes by \(0,-1\), or
\(-2\), accounting for the preceding zombie move. A gap at most \(1\)
results in immediate or next-turn capture. Hence, while evasion
continues, \[
2\le g_t\le p\le d+1.
\qquad\text{(3)}
\] These bounds also ensure that the first zombie's move in the chosen
direction is always its unique shortest-path move.

For \(0\le t\le d\), (3) gives \[
0<u_t\le t+d+1\le2d+1<2d+3.
\] Thus the survivor cannot reach the attachment vertex during these
rounds and cannot leave the cycle.

The outside zombie consequently reaches the root after exactly \(d\)
moves. After the survivor's response on that round, \[
d+2\le u_d\le2d+1.
\] The zombies are now at coordinates \(d\) and \(0\). They bracket the
survivor on the arc from \(d\) forward to \(0\), and the two distances
along that arc are \[
u_d-d\le d+1,
\qquad
2d+3-u_d\le d+1.
\] Both are shortest-path distances.

The zombies now move inward along this arc. Each of these two distances
is nonincreasing after a complete round, so the inward moves remain
legal. The occupied endpoints prevent escape through the root, and the
intervening arc loses two edges on every zombie turn. Capture follows.
\(\square\)

\begin{center}\rule{0.5\linewidth}{0.5pt}\end{center}

\subsubsection{\texorpdfstring{5. Two zombies win on
\(H\)}{5. Two zombies win on H}}\label{two-zombies-win-on-h}

Place the zombies at \[
v_0\quad\text{and}\quad v_{11}.
\]

Project every vertex of an attached cycle to its attachment vertex on
the path \[
P=v_0v_1\cdots v_{11}.
\]

Until a zombie reaches the attachment vertex of the cycle currently
containing the survivor, both zombies move along \(P\) toward the
survivor's projection. These are legal shortest-path moves.

The survivor cannot cross an occupied path vertex. Consequently, unless
capture or such a cycle-root contact occurs, the two zombies' positions
after \(t\) moves are \[
v_t,\qquad v_{11-t},
\] and the survivor's projection remains strictly between them. This
shrinking interval cannot support evasion indefinitely.

It remains to handle the first cycle-root contact. All attachment
indices are \(0,2,4\), so this first contact is made by the zombie
coming from \(v_0\). If the relevant root is \(v_r\), the contact occurs
after \(r\) zombie moves, with positions \[
v_r,\qquad v_{11-r}.
\] The distance of the second zombie from the root is \[
d=11-2r.
\]

The three possibilities are \[
\begin{array}{c|c|c}
r&d&\text{Length of the attached cycle}\\ \hline
0&11&25=2\cdot11+3\\
2&7&17=2\cdot7+3\\
4&3&9=2\cdot3+3
\end{array}
\]

If the survivor is actually at the root, it is captured. Otherwise it
lies inside that cycle and cannot leave through the occupied root on its
ensuing response. The capture lemma therefore applies in every case.

Thus \[
z(H)\le2.
\]

Conversely, one zombie can always be evaded on \(A\): in the rotation
lemma, a single phase can always be accommodated by choosing
\(q=p+2\pmod{25}\). This applies whether the zombie starts inside or
outside \(A\). Hence \[
\boxed{z(H)=2.}
\]

\begin{center}\rule{0.5\linewidth}{0.5pt}\end{center}

\subsubsection{6. Conclusion}\label{conclusion}

The explicitly constructed graphs satisfy \[
\boxed{z(H)=2<3\le z(G),}
\] although \(H\) is obtained from \(G\) by adding just the leaf
\(v_{11}\) adjacent to \(v_{10}\).

Thus \textbf{zombie number is not invariant under leaf attachment}, even
for cactus graphs.

The mechanism is precisely an initial-delay advantage: the new leaf
delays the rear zombie by one move, making its arrival gap at each of
the three cycle roots large enough to approach the survivor from the
opposite direction. The lower-bound argument shows that no alternative
pair of starting vertices in \(G\) can achieve a global win.

Every pair of starting positions in \(G\) and every survivor start in
\(H\) is covered above. Determining the exact value of \(z(G)\) beyond
\(z(G)\ge3\) is unnecessary for the counterexample.

\appendix

\section{Provenance and review
status}\label{provenance-and-review-status}

\begin{warning}

\textbf{UNREFEREED MODEL OUTPUT.} This document was selected solely
because GPT-6 Astra labelled its own result
\texttt{would\_publish:\ true} and returned \texttt{proved} or
\texttt{disproved}. It has not passed the adversarial LLM referee used
for the earlier Sol campaign, has not been checked by a human
mathematician, and has not been checked for novelty. Treat every
mathematical and bibliographic claim as unverified.

\end{warning}

{\def\LTcaptype{none} % do not increment counter
\begin{longtable}[]{@{}
  >{\raggedright\arraybackslash}p{(\linewidth - 2\tabcolsep) * \real{0.5000}}
  >{\raggedright\arraybackslash}p{(\linewidth - 2\tabcolsep) * \real{0.5000}}@{}}
\toprule\noalign{}
\endhead
\bottomrule\noalign{}
\endlastfoot
Catalog id & \texttt{2008.03587\_\_00} \\
Catalog entry &
\href{https://graph-theory-ai.github.io/graph-conjectures/arxiv/2008.03587__00/}{Question
4.1} \\
Source corpus & arxiv \\
Campaign leg & selected second attempts (\texttt{attacks\_retry}) \\
Source paper / entry & A note on deterministic zombies \\
Model verdict & \textbf{disproved} (confidence: high) \\
Selection criterion & \texttt{would\_publish:\ true} and verdict
\texttt{proved} or \texttt{disproved} \\
Model & \texttt{gpt-6-astra}, reasoning effort \texttt{max},
\texttt{mode=pro}, flex service tier \\
Original artifact &
\texttt{attacks\_retry/2008.03587\_\_00/output.md} \\
Independent review & \textbf{None} \\
arXiv & \href{https://arxiv.org/abs/2008.03587}{arXiv:2008.03587} \\
Previous attempt & \texttt{attacks/2008.03587\_\_00}: gpt-5.6-sol
verdict \texttt{partial} \\
Model's caveats & Uses the standard compulsory-move rules, with zombies
moving before the survivor; no computational or literature-status claim
is made. \\
\end{longtable}
}

The generated Markdown and LaTeX sources for this PDF are stored under
\texttt{to\_review\_astra/src/2008.03587\_\_00/}.

\end{document}
